
\capitulo{Proceso Ingeniería de Requerimientos}
\label{sec:srs-4}

El proceso de ingeniería de requerimientos se fundamenta en el estándar de ingeniería de requerimientos~\cite{iso29148_2018} y se desarrolla de manera iterativa, conforme a la metodología Disciplined Agile Delivery (DAD). Comprende la identificación de necesidades, especificación, verificación, validación y gestión de los requerimientos del producto.

\subsection{Análisis de la misión}
\label{sec:srs-4-1}

El sistema busca apoyar la comprensión del posible comportamiento de incendios en el sector de Aguas Claras de los Cerros Orientales de Bogotá, mediante una aplicación de escritorio que consulte escenarios precalculados con FDS y presente estimaciones georreferenciadas. Sus resultados están sujetos a las limitaciones de las simulaciones y no sustituyen el criterio profesional de los responsables de la emergencia.

\subsection{Necesidades de los interesados}
\label{sec:srs-4-2}

Las necesidades se identificaron mediante el análisis del problema, la revisión de antecedentes y entrevistas con profesionales relacionados con la atención de emergencias. El usuario principal es el bombero coordinador, quien requiere consultar condiciones meteorológicas, localizar incendios y visualizar estimaciones de propagación sin conocimientos especializados en simulación. Estas necesidades se revisarán mediante evaluaciones con usuarios.

\subsection{Requisitos del sistema y del software}
\label{sec:srs-4-3}

Las necesidades identificadas se transforman en requerimientos funcionales y no funcionales, documentados en la sección 3. Se procura que cada requerimiento sea identificable, claro, consistente y verificable, siguiendo los principios de la norma citada y considerando las restricciones técnicas, los atributos de calidad y las funcionalidades previstas.

\subsection{Verificación y validación}
\label{sec:srs-4-4}

La verificación comprobará la claridad, consistencia y trazabilidad de los requerimientos mediante revisiones de la especificación. La validación evaluará su correspondencia con las necesidades de los usuarios, incluyendo las pruebas de usabilidad con bomberos previstas en la sección 3.1.1. Los resultados permitirán identificar ajustes necesarios.

\textcolor{black}{Según el cronograma del PMP, la evaluación con usuarios se realiza en I6 sobre la versión integrada y \textcolor{black}{el dataset congelado}. Las observaciones se revisan el 23 de abril de 2027. En I7 se verifican las correcciones y se consolidan la tabla de referencia cruzada y los registros de pruebas para el cierre con el director el 7 de mayo. La entrega académica del paquete corresponde a la semana 17.}

\subsection{Gestión de requisitos}
\label{sec:srs-4-5}

Los requerimientos se gestionarán durante las iteraciones del proyecto mediante control de versiones y una tabla de referencia cruzada que relacione necesidades, funcionalidades, casos de uso y pruebas. Los cambios se evaluarán según su impacto en el alcance y las dependencias técnicas, manteniendo actualizada la documentación correspondiente.
